\chapter{Conclusion}
\label{chap:conclusion}

%% PURPOSE: This chapter closes the dissertation by restating what was actually
%% achieved (contributions) and what remains open (future work). It should be
%% concise and free of hedging — by this point the reader has seen the evidence.

\section{Summary of Contributions}
\label{sec:contributions}
% TODO: Enumerate the concrete, citable contributions of the dissertation.
% Be specific — avoid vague claims like "this thesis advances the field."
% Suggested structure (adapt based on the final results):
%
%   (1) A validated, open-source pipeline for discovering virtual choreographies
%       from K-12 online learning event logs, demonstrated to recover known
%       behavioral archetypes from synthetic data with ARI ~ 0.53 and NMI ~ 0.59.
%
%   (2) Identification and characterization of four virtual choreographies in the
%       synthetic Connections Academy dataset (Steady_Worker, Balanced_Engager,
%       Low_Engagement, Night_Crammer), with quantified outcome associations and
%       statistical effect sizes.
%
%   (3) Evidence that temporal regularity and synchronous-participation balance
%       are positively associated with academic success (supporting H1, H2),
%       while night-heavy and low-engagement patterns correlate with weaker
%       outcomes (supporting H5).
%
%   (4) A first-N-weeks early prediction study showing that choreography labels
%       add predictive value (ROC-AUC gain) over volume metrics alone as early
%       as the first four weeks of enrolment (addressing RQ4).
%
%   (5) A multi-method interpretability analysis (permutation importance + SHAP)
%       that identifies the specific behavioral signals most actionable for
%       on-site teacher facilitators, directly addressing the actionability
%       requirement of RQ5.
%
%   (6) An honest characterization of the scope of the synthetic study, including
%       the conditions under which the findings can be expected to transfer to
%       the real Pearson log data.

\section{Future Work}
\label{sec:future-work}
% TODO: Describe the planned and recommended next steps. Order from most
% immediate to most exploratory:
%
%   (1) Application to real Pearson / Connections Academy log data — the primary
%       planned follow-up. This will validate whether the pipeline generalizes
%       beyond the synthetic setting and produce empirical claims about actual
%       K-12 students.
%
%   (2) Longitudinal choreography tracking — follow students' cluster membership
%       across multiple semesters or academic years. Do choreographies shift after
%       interventions? Can transitions (e.g., Low_Engagement → Balanced_Engager)
%       be predicted early enough to trigger support?
%
%   (3) Richer feature engineering — incorporate interaction quality features
%       beyond counts: response latency to teacher feedback, forum post length,
%       revision cycles on submitted assignments, and synchronous attendance rate.
%       These may improve cluster separation and outcome prediction.
%
%   (4) Teacher-facing early-warning dashboard — translate the interpretability
%       results into actionable, student-level risk scores for on-site facilitators,
%       with choreography-specific intervention recommendations.
%
%   (5) Subject-level choreography analysis — re-run the pipeline within individual
%       subjects (e.g., Mathematics vs. English Language Arts) to test
%       H\_Robustness more rigorously and identify subject-specific behavioral
%       signatures.
%
%   (6) Comparison with supervised approaches — once human-validated choreography
%       labels are available from the real data, compare unsupervised discovery
%       against supervised multi-label classification to quantify how much is
%       gained or lost by requiring labeled training data.
